\documentclass[11pt]{article}
\usepackage[margin=1in]{geometry}
\usepackage{amsmath,amssymb,booktabs,graphicx,microtype}
\usepackage[colorlinks=true,linkcolor=blue!50!black,citecolor=blue!50!black]{hyperref}
\usepackage{caption}
\captionsetup{font=small}
\title{\textbf{Prediction and inference in ProductGPT}\\[2pt]
\large An executive report on what the model can claim}
\author{ProductGPT technical note}
\date{29 September 2026}

\begin{document}
\maketitle

\begin{abstract}
\noindent ProductGPT has been developed and selected, for a year, on one
criterion: how well it predicts the \emph{next} decision given a true history.
That programme is now finished, and it has converged. This report separates
what that work established (the \textbf{prediction phase}) from what a
marketing paper actually needs (the \textbf{inference phase}), and reports the
evidence for a claim that has emerged repeatedly and from independent
directions: \emph{the criterion that selected the model cannot select a model
for inference}. Four experiments completed on 29 September make this a
measurement rather than an argument. The infrastructure to test the remaining
step --- whether the model is valid as a \emph{generator} of sequences --- is
built and validated, and has not yet been run.
\end{abstract}

\section{The one-paragraph version}

Every architecture we tried reaches the same predictive accuracy. Every
component we added to improve it --- the inventory mechanism, customer
heterogeneity, the learned product embedding --- is worth nothing or nearly
nothing on that criterion. Yet the same components are decisive for whether the
model's parameters mean anything: modelling customer heterogeneity moves the
recovery of the substitution structure from essentially zero to clearly found,
and removes a factor-2.2 bias in the satiation coefficient. Selecting on
predictive fit would have discarded the single component the identification
analysis says matters most. That is the finding, and it reframes the paper:
the identification analysis is the contribution, and the architecture is the
setting in which it is made.

\section{The prediction phase: what was established}

\paragraph{The headline is a tie, and it is a strong one.} Four model families
--- a plain GRU, a GRU encoder, a transformer with recency bias, and a hybrid
--- were each tuned to their own frontier under an equal budget, then evaluated
once on a holdout of customers and campaigns neither had seen. All eight cells
(four families $\times$ two vocabularies) land within 0.0194 nats, against a
seed standard deviation of 0.0075. No family wins. The best cell is also the
smallest: 1.25M parameters against 7.93M for the transformer.

\paragraph{The components do not separate either.} Removing the stock path
entirely --- the inventory representation that three rounds of work went into
--- costs 0.0137 nats, inside the same margin.

\paragraph{The four experiments completed on 29 September.} Thirty GPU runs and
one CPU job, five seeds per cell, with selection read on validation and the
holdout opened once and not used to choose.

\begin{center}\small
\begin{tabular}{llr}
\toprule
Experiment & Question & Validation NLL \\
\midrule
control & frozen configuration, half-life 2 & $0.8756 \pm 0.0085$ \\
batch 30 & mixture head over output projections & $0.8756 \pm 0.0097$ \\
batch 30 & per-customer embedding & $0.8802 \pm 0.0076$ \\
batch 30 & both & $0.8829 \pm 0.0116$ \\
batch 32 & attributes only, no product identity & $0.8757 \pm 0.0078$ \\
batch 26 & product identity only, no attributes & $0.8820 \pm 0.0101$ \\
\bottomrule
\end{tabular}
\end{center}

\noindent Three readings, all negative for prediction:

\begin{enumerate}
\item \textbf{Customer heterogeneity buys nothing.} Three variants, five seeds
each; the mixture head ties the control exactly and the other two are worse.
Every difference is inside a seed standard deviation, but the direction is
uniformly non-positive and no variant is adopted.

\item \textbf{The learned per-product embedding buys nothing.} Deleting it
costs $0.0001$ nats and is slightly better on the holdout ($0.8832$ against
$0.8846$); deleting the product \emph{attributes} instead costs $0.0064$.
Everything the model knows about a product, it knows from the attribute
projection. This confirms by construction an earlier decomposition in which the
learned identity branch scored at chance, and it retires a product map we had
previously shown as evidence of learned structure.

\item \textbf{The one positive result is a boundary check.} A profile over
frozen decay half-lives gives $0.8795$, $\mathbf{0.8735}$, $0.8756$, $0.8801$
and $0.8817$ at half-lives of $0.5$, $1$, $2$, $5$ and $10$ occasions. The
optimum is \emph{interior}: holdings fade over days, not campaigns. The
$1$-versus-$2$ gap is a quarter of a seed standard deviation, so the
defensible claim is a half-life of one to two occasions, not a point estimate.
\end{enumerate}

\paragraph{What the prediction phase licenses.} Forecasting the next decision.
That is the whole of it. It is a real result --- the ceiling claim is
pre-registered and now supported at five seeds after tuning --- but it supports
no statement about why a customer behaves as they do, and no policy number.

\section{The inference phase: three kinds of validity}

The word ``valid'' has been doing too much work. Three distinct claims, with
distinct evidence and distinct licences:

\begin{center}\small
\begin{tabular}{p{2.6cm}p{4.3cm}p{4.3cm}}
\toprule
Tier & Evidence & What it licenses \\
\midrule
T1 decision-level predictive & stage 4/5: NLL, AUC, calibration error
0.010--0.019 & forecasting the next decision \\
T2 sequence-level generative & \textbf{not yet tested} & in-support policy
simulation: pacing, campaign length, re-ordering assortments that occurred \\
T3 structural & decay and duplicate weights recovered; a free product-level
kernel is not & off-support counterfactuals: new pairings, novel assortments \\
\bottomrule
\end{tabular}
\end{center}

\noindent T2 is neither implied by T1 nor dependent on T3. Every result we hold
is T1 or T3. The middle rung --- the one that licenses the policy numbers a
marketing paper wants --- has never been tested.

\paragraph{Why T1 does not imply T2.} If the model's conditional distributions
were correctly specified, the chain rule would deliver the joint law for free.
The gap opens only under misspecification, and for a specific reason: training
evaluates the model on prefixes drawn from the \emph{data}, while a rollout
visits prefixes drawn from the \emph{model}. A sequence-level failure in the
presence of good one-step fit is therefore direct evidence of misspecification
--- a specification test that the likelihood cannot perform, because it is
never evaluated off the data manifold. Marketing has long had the same idea
under other names: nonrandom-holdout validation, predictive simulation for
discrete choice models, and posterior predictive checks.

\section{What the identification work establishes (T3)}

\paragraph{Some parameters are recoverable and some are not.} In simulations
where the truth is known by construction, the decay and the duplicate weights
are recovered; a free product-level substitution kernel is not, and a rank
condition on the offer calendar says it never will be. Our window offers 30
campaigns against 49 products; the full six-year history offers 106 against
135. \textbf{The deficit widens as the panel grows}, because a live-service
game introduces products about as fast as it introduces campaigns. Waiting does
not help, and neither does an independently constructed banner schedule, which
reproduces the deficit exactly.

\paragraph{Heterogeneity, not the calendar, is the binding constraint.} The
most consequential simulation result is not about assortments at all. With
unobserved taste heterogeneity that the estimator does not model, the recovered
attribute structure collapses:

\begin{center}\small
\begin{tabular}{lrrr}
\toprule
Unobserved taste dispersion & 0.25 & 0.50 & 1.00 \\
\midrule
attribute lift, tastes not modelled & 0.15 & 0.03 & 0.18 \\
attribute lift, tastes modelled & 28.83 & 27.85 & 10.56 \\
satiation coefficient, not modelled (true $-1.50$) & $-1.28$ & $-1.35$ &
$\mathbf{-3.35}$ \\
satiation coefficient, modelled & $-1.31$ & $-1.30$ & $-1.28$ \\
\bottomrule
\end{tabular}
\end{center}

\noindent Modelling per-household attribute tastes restores the structure and
removes the factor-2.2 bias in the satiation coefficient. One caveat belongs
with these numbers: the true lift is 7.4, and the recovered values overshoot
it. The qualitative claim is solid; the magnitudes are noisy at three seeds and
should not be quoted as estimates.

\section{The two phases meet}

\begin{center}\small
\begin{tabular}{p{2.8cm}p{5.0cm}p{5.0cm}}
\toprule
Component & Predictive verdict & Structural verdict \\
\midrule
Customer heterogeneity & No gain. Five seeds, three variants, none beats the
control. & Decisive. Without it the attribute kernel is unrecoverable and
satiation is biased by 2.2$\times$. \\
Product representation & Identity and attributes indistinguishable; deleting
identity costs 0.0001. & Only the attribute parameterisation is identified
under our rotation. \\
\bottomrule
\end{tabular}
\end{center}

\noindent Two independent cases --- one from simulation, one from the real
panel --- of decision-level accuracy being silent about exactly the thing the
paper wants to claim. Neither was constructed to make the point; both arrived
from experiments run for other reasons, which is what makes them usable.

\section{What is built, and what remains}

\paragraph{Built and validated this week, not yet run.} A rollout engine that
carries a trained model forward on its own outputs across the holdout
campaigns; an acquisition environment that turns a decision into what the
player received, with the lottery dynamics taken from the published rates and
the product-identity process calibrated on calibration data only; twenty-seven
sequence-level statistics the model was never fitted to, including the
satiation signature --- the probability of buying on a banner given how many
copies of its featured product the customer already holds; and a null band
built by splitting the real data against itself, without which no discrepancy
number means anything.

The instrument is calibrated: it rejects an independent-draws generator
overwhelmingly, while a first-order Markov chain scores less than half that
discrepancy and still fails. A 9$\times$9 transition matrix is therefore a live
benchmark, not a straw man.

\paragraph{The test that matters.} Four model families tie on one-step accuracy.
Do they tie as generators? The satiation mechanism is a sequence-level
mechanism --- it governs whether a customer who has just acquired five copies
stops --- so it is precisely the kind of thing a decision-level criterion is
least able to see. If tied models separate here, the criterion was wrong rather
than the mechanism. If they do not, the case for using the model as a simulator
fails cleanly, and that is also a result.

\paragraph{One scope decision, which gets more expensive to defer.} The model
consumes elapsed time as an input and never predicts it, so it cannot generate
the timing of occasions. Everything built so far conditions on the observed
arrival process, and therefore tests the composition and dependence structure
of decisions rather than purchase timing or frequency. Removing that limitation
requires adding a duration head --- a modelling change, not an evaluation
change, and one that would make ProductGPT a genuine sequence generator.

\section{What this means for the paper}

The spine that all of this supports is: \emph{a deep sequence model of
limited-time product markets, and an account of what such a model can and
cannot learn from observational choice data}. This inverts the usual order ---
the identification analysis is the contribution and the architecture is the
setting --- and every result, including the negative ones, supports it.

The claim to avoid is predictive superiority. Our own stage-4 tie forecloses
it, and it is the crowded lane. A tie is fatal to that spine and irrelevant to
this one.

\end{document}
